\documentclass[10pt]{article}
\providecommand{\FIGDIR}{fig}
\newcommand{\DIGESTDATE}{5 November 2026}
\input{preamble}

\begin{document}

% ==================================================================== MASTHEAD
\thispagestyle{empty}
\begin{tikzpicture}[remember picture, overlay]
  \fill[Deep] (current page.north west) rectangle
      ($(current page.north east)+(0,-67mm)$);
  \fill[Amber] ($(current page.north west)+(0,-67mm)$) rectangle
      ($(current page.north east)+(0,-68.4mm)$);
  \foreach \x/\y/\r in {22/12/0.9, 41/26/1.5, 63/9/0.7, 88/31/1.1, 112/17/1.9,
                        138/29/0.8, 159/13/1.3, 178/54/1.0, 196/21/1.6,
                        30/44/0.6, 74/57/1.2, 125/48/0.9, 168/6/0.8, 191/59/0.7,
                        52/18/1.0, 100/60/0.6, 148/20/1.1, 66/49/1.4, 205/45/0.8}
    \fill[Sky, opacity=0.30] ($(current page.north west)+(\x mm,-\y mm)$) circle (\r mm);
\end{tikzpicture}

\vspace*{4mm}
{\sffamily\fontsize{7.6}{9}\selectfont\color{Sky}%
 AUTOMATED DAILY LITERATURE DIGEST\;\textbullet\;PhD RESEARCH WATCH\par}
\vspace{2.2mm}
{\sffamily\bfseries\fontsize{27}{29}\selectfont\color{white}%
 Particulate Matter\par}
{\sffamily\fontsize{27}{29}\selectfont\color{Sky}%
 Monitoring \& Health Effects\par}

\vspace{5.0mm}
{\sffamily\fontsize{8.4}{10}\selectfont\color{white}%
 Issue 2026-11-05\;\;\textcolor{Amber}{\textbar}\;\;Entry window 5 November 2026 (full day)\;\;%
 \textcolor{Amber}{\textbar}\;\;25 records screened in scope\par}
\vspace{18mm}

\noindent
\begin{minipage}[t]{0.190\linewidth}\metric{25}{RECORDS IN SCOPE\\(63 EXCLUDED)}{Deep}\end{minipage}\hfill
\begin{minipage}[t]{0.190\linewidth}\metric{9}{OF 10 SUBTOPIC\\BANDS POPULATED}{Teal}\end{minipage}\hfill
\begin{minipage}[t]{0.190\linewidth}\metric{6}{EFFECT ESTIMATES\\WITH A CI}{Sky}\end{minipage}\hfill
\begin{minipage}[t]{0.190\linewidth}\metric{2}{TIER-A\\RECORDS}{Amber}\end{minipage}\hfill
\begin{minipage}[t]{0.190\linewidth}\metric{5}{RECORDS CARRIED ON\\METADATA ALONE}{Clay}\end{minipage}

\vspace{6mm}

\section*{\sffamily\bfseries\color{Deep}Signal of the day}
\vspace{-1.5mm}{\color{Amber}\rule{\linewidth}{1.1pt}}\vspace{2.5mm}

\begin{keybox}
{\sffamily\bfseries\small\color{Deep}A dense low-cost PM$_{2.5}$ network, gap-filled by machine learning, was carried all the way to district-level short-term mortality across Bihar in \textit{Nature Communications}: northern districts ran 1.2 times the mean exposure of southern ones, district mortality varied 1.7-fold and daily attributable deaths 3.2-fold.}\\[1.4mm]
\footnotesize
IIT Kanpur (Tripathi group); network PM$_{2.5}$ fused with NFHS socioeconomic data and India-specific relative risks. State-level exposure inequity was limited; inter- and intra-regional inequity within Bihar was substantial.
\textbf{Where it is weak:} the abstract reports no node count, no calibration R$^2$ or RMSE and no uncertainty on attributable deaths. The mortality is a risk function applied to a sensor-derived surface, not an observed association, so any network calibration bias passes one-for-one into the burden.
\textbf{Why it matters here:} this is the clearest high-visibility case of low-cost sensing as the exposure layer of a burden estimate outside cities. Its value to the field rests on the calibration and propagated-uncertainty detail in the full text, which the weekly should pull.\par
\end{keybox}

\vspace{3mm}

\section*{\sffamily\bfseries\color{Deep}What changed today}
\vspace{-1.5mm}{\color{Amber}\rule{\linewidth}{1.1pt}}\vspace{2.5mm}

\footnotesize
\begin{itemize}

\item \textbf{Sinonasal cancer, and a window gradient.} US commercial claims, 3,600 cases and 17,695 controls: per SD of PM$_{2.5}$ (about 2.5\,$\mu$g/m$^3$), aOR 1.25 (1.05--1.49) at 6 months, 1.31 (1.08--1.59) at 1 year, 1.36 (1.11--1.66) at 5 years. Exposure is a 1-km satellite-fused product collapsed to metropolitan area, case-control contrast is ~0.2--0.3\,$\mu$g/m$^3$, and wood and leather dust, the main SNM cause, is unmeasured.

\item \textbf{Two influenza designs, one PM magnitude.} Shanxi 11-city time-series: RR 1.024 (1.015--1.034) per 10\,$\mu$g/m$^3$ PM$_{2.5}$, but NO$_2$ (1.115) and SO$_2$ drive the mixture. Shanghai laboratory-confirmed case-crossover (\textsc{backfill}): PM$_{10}$ OR 1.019 (1.004--1.034). Both near 2\,\% per 10\,$\mu$g/m$^3$, both on regulatory monitors; gases carry more signal than PM.

\item \textbf{Where the dose is, mass is not.} Chamber emission of six indoor sources: most mass and the highest mass-normalised DTT OP sit below 0.25\,$\mu$m, and for equal PM$_{10}$ mass these sources give an estimated 76 $\pm$ 10\,\% more alveolar dose than ambient PM. Kanpur--Mumbai impactor pairs: fine mode delivers 53--75\,\% of alveolar dose from 40--62\,\% of mass. Both point at the size range optical sensors count worst.

\item \textbf{Composition breaks transfer.} Around an opencast coal mine (\textsc{backfill}): roadside PM$_{2.5}$ near 400\,$\mu$g/m$^3$, Cr, Ni, Fe, Zn 5--10-fold above background, OP only weakly explained (25--40\,\%). Shenzhen fireworks: 4.7\,\% of PM$_{2.5}$ over the season but $>$50\,\% at peak hours. Accra FRM series (\textsc{backfill}): three sites at 3--4 times the WHO guideline, a reference anchor for West African sensor calibration.

\item \textbf{Exposure basis (x1).} Of 25 records, 1 rests on a low-cost network (Bihar, on the sensing side of x1) and none on a personal monitor. Health side by subtopic (n=7 with a recoverable basis): 71\,\% measured, 14\,\% satellite-fused, 14\,\% proxy --- but the measured share is three regulatory-monitor respiratory studies and two occupational filter campaigns. \emph{All six effect estimates below come from regulatory monitors or an MSA-averaged satellite product}; none is tied to a sensor-measured exposure.

\end{itemize}

\vspace{2mm}

\begin{keybox}
{\sffamily\bfseries\footnotesize\color{Deep}Window continuity}\\[1.2mm]
{\fontsize{7.2}{8.8}\selectfont
The 4 November issue closed with \texttt{last\_entry\_date} at 4 November. This issue was built on the \textbf{6 November 17:34 UTC (11:34 CST) run}, the 05:04 UTC slot firing 12.5 hours late. Before 22:00 the newest buildable day is yesterday, so \textbf{5 November is covered in full}. Thursday: no rollup is owed; \textbf{W45 (1--7 November) falls on Saturday 7 November}. Three \textsc{backfill} records, the cap: two held from the 4 November run (Accra FRM, Shanghai influenza) and one from Amass (coal-mine PM$_{2.5}$ composition). \textbf{Europe PMC returned zero records for the day} (likely indexing lag; the leg is not broken on 4 Nov). The OpenAlex connector was again not available on this session.\par}
\end{keybox}

\vspace{3mm}
\par\ifdim\dimexpr\pagegoal-\pagetotal\relax<18\baselineskip\newpage\fi
\section*{\sffamily\bfseries\color{Deep}Corpus analytics}
\vspace{-1.5mm}{\color{Amber}\rule{\linewidth}{1.1pt}}\vspace{2.5mm}

\noindent\begin{minipage}{\linewidth}\includegraphics[width=\linewidth]{\FIGDIR/f1_subtopics.png}
\figcap{Subtopic assignment of the 25 records across nine of ten bands: exposure assessment 9, sensing 4, respiratory 3, occupational and indoor 3, burden 2, and one each for other clinical, neuro, mechanistic and cardiovascular. Reproductive is empty today.}\end{minipage}

\vspace{2mm}
\noindent\begin{minipage}{\linewidth}\includegraphics[width=\linewidth]{\FIGDIR/f2_design_endpoint.png}
\figcap{Left --- architecture: measurement campaign 6, review or synthesis 6, metadata-only 4, and 2 each cross-sectional, acute observational, modelling and experimental; one cohort (CAPNETZ). Right --- endpoints are singletons apart from 5 unrecoverable (four Elsevier metadata-only records and the UFP letter); nine of the twenty labels are measurement or chemistry outcomes, not clinical ones.}\end{minipage}

\vspace{2mm}
\noindent\begin{minipage}{\linewidth}\includegraphics[width=\linewidth]{\FIGDIR/f3_metric_geography.png}
\figcap{Left --- particle metric: PM$_{2.5}$ only 10, unspeciated or emissions 6, PM$_{2.5}$+PM$_{10}$ jointly 5, size-resolved 2, optical 1, ultrafine 1. Right --- geography: global or multi-region 7 (four reviews, the UFP letter, one metadata-only record with no location and the laboratory AMS study, which fall there by default), South Asia 5 (India 4, Pakistan 1 counting the regional review), China 5, North America 4, sub-Saharan Africa 2, and one each Europe and Southeast Asia.}\end{minipage}

\vspace{2mm}
\noindent\begin{minipage}{\linewidth}\includegraphics[width=\linewidth]{\FIGDIR/f6_lifecourse.png}
\figcap{Life-course windows. The two influenza series are all-age and are counted in every postnatal stage; working age adds firefighters, the dust review and the sinonasal case-control (case mean age 51); older adults add the CAP cohort and the Alzheimer review. No preconception or in-utero record today.}\end{minipage}

\vspace{2mm}
\noindent\begin{minipage}{\linewidth}\includegraphics[width=\linewidth]{\FIGDIR/f4_heatmap.png}
\figcap{Subtopic $\times$ study architecture. Exposure assessment splits across measurement (3), metadata-only (3), modelling (2) and the letter; sensing is two measurement records, one metadata-only and the AMS calibration. The only cohort is respiratory (CAPNETZ), and neither respiratory acute design uses a sensor.}\end{minipage}

\vspace{2mm}

\IfFileExists{\FIGDIR/extra_figs.tex}{\input{\FIGDIR/extra_figs.tex}}{}



\noindent\begin{minipage}{\linewidth}\includegraphics[width=\linewidth]{\FIGDIR/f5_forest.png}
\figcap{Ratio scale, 6 estimates from 3 sources: Yang sinonasal case-control (3, per SD), Zhang Shanxi influenza (2, per 10\,$\mu$g/m$^3$), Hao Shanghai influenza (1, per 10\,$\mu$g/m$^3$). The three sinonasal estimates are nested averaging windows of one exposure in one sample and are not independent. The three influenza estimates share endpoint family and increment, but two are PM$_{10}$ and one PM$_{2.5}$, and RR from a time-series is not OR from a case-crossover; no pooling and no I$^2$.}\end{minipage}

\vspace{4mm}

\par\vskip 0pt plus 24\baselineskip\penalty-200\vskip 0pt plus -24\baselineskip
\section*{\sffamily\bfseries\color{Deep}Paper-level digest}
\vspace{-1.5mm}{\color{Amber}\rule{\linewidth}{1.1pt}}\vspace{2.5mm}

{\fontsize{7.4}{9}\selectfont\color{Slate}Ordered by cluster. Each entry gives the
methodological core and the finding that matters, not an abstract paraphrase. Records
marked \textit{metadata only} carry no recoverable abstract and assert no finding.\par}

\band{Neuro / mental health\hfill 1 record}{Deep}

\paperentry{Defrancesco lung-brain axis review PM$_{2.5}$ Alzheimer disease mechanisms 2026}
{Mechanistic understanding of the role of the lung-brain axis in the pathogenesis of PM$_{2.5}$-induced Alzheimer's disease.}
{Clin Sci (Lond)}
{Narrative review positioning pulmonary injury, rather than direct translocation, as the main route from inhaled PM$_{2.5}$ to AD pathology: systemic inflammation, oxidative stress, blood-brain barrier disruption and extracellular-vesicle signalling converge on amyloid, tau, microglial and autophagy-lysosomal dysfunction. Modifiers named are composition, co-pollutants, sex, APOE and epigenetic reprogramming. No search protocol, no quantitative synthesis; the framework is a hypothesis map and the human causal link remains epidemiological association. Relevance: indirect - if the pulmonary route dominates, deposited dose in the alveolar region, not ambient mass, is the exposure metric a sensor should be calibrated to serve.}
{10.1042/cs20262109}

\par\vskip 0pt plus 10\baselineskip\penalty-200\vskip 0pt plus -10\baselineskip
\band{Cardiovascular \& metabolic\hfill 1 record}{Teal}

\paperentry{Siddiqui narrative review AQI ASCVD South Asians 2026}
{Impact of Air Quality Index on Atherosclerotic Cardiovascular Disease Risk in South Asians: A Narrative Review.}
{Curr Atheroscler Rep}
{Narrative review arguing that worse AQI and its constituent pollutants raise ASCVD risk in South Asians, who already carry high genetic and metabolic predisposition, and that pollution may compound traditional risk factors. Reports no search strategy, no pooled estimate and no South Asia-specific effect sizes; it concedes that region-specific intervention evidence is thin. Using AQI as the exposure is the core weakness: it is a max-of-pollutants communication index, not a dose metric, and conflates PM with gases. Relevance: indirect - the region with the highest burden still lacks cohort-grade exposure, which is where dense sensing would add most.}
{10.1007/s11883-026-01486-0}

\par\vskip 0pt plus 10\baselineskip\penalty-200\vskip 0pt plus -10\baselineskip
\band{Respiratory \& allergic\hfill 3 records}{Sky}

\paperentry{Zhang Shanxi 11-city time-series pollutant mixtures influenza vaccination 2026}
{Influenza Incidence in Relation to Air Pollution Mixtures and the Modifying Role of Vaccination in Shanxi Province, China.}
{Environ Res}
{11 Shanxi cities, 2014-2022, DLNM with zero-inflated negative binomial and meta-analysis. Per 10 $\mu$g\,m$^{-3}$ PM$_{2.5}$, RR 1.024 (1.015-1.034); PM$_{10}$ 1.018 (1.010-1.025); NO$_2$ 1.115 (1.086-1.146). WQS (1.078, 1.005-1.157) and qgcomp (1.122, 1.064-1.182) mixture effects were driven by SO$_2$ and NO$_2$, not PM. Vaccination modification is city-level coverage, therefore ecological, and labelled exploratory. Surveillance incidence is care-seeking and testing dependent; 2020-2022 includes COVID-era behavioural shifts the abstract does not say were handled. Relevance: indirect - PM signal is small next to gases, which argues for multi-pollutant nodes rather than PM-only networks for infectious endpoints.}
{10.1016/j.envres.2026.126109}

\paperentry{Hao Shanghai case-crossover lab-confirmed influenza six pollutants 2026}
{Association between short-term air pollutant exposure and influenza incidence in Shanghai, China: A case-crossover study}
{Int J Biometeorol}
{Laboratory-confirmed influenza cases, Shanghai 2013-2017, time-stratified case-crossover with IDW exposure at the individual level. Per 10 $\mu$g\,m$^{-3}$ PM$_{10}$, OR 1.019 (1.004-1.034); O$_3$ was inverse (OR 0.959, 0.941-0.977). PM$_{2.5}$, SO$_2$, NO$_2$ and CO were positive across lags but their point estimates are not in the abstract. Laboratory confirmation is the design's strength over syndromic counts; IDW from a sparse regulatory grid and selection into testing are the weaknesses, and the negative O$_3$ term likely absorbs season. Read alongside the same-day Shanxi series: two independent Chinese designs give PM effects near 2\% per 10 $\mu$g\,m$^{-3}$. Relevance: indirect. \textsc{backfill}, entry 20 March 2026.}
{10.1007/s00484-026-03163-y}

\paperentry{Kutzinski CAPNETZ air pollution sputum pathogen spectrum community-acquired pneumonia 2026}
{Air pollution is associated with a distinct pathogen spectrum in community-acquired pneumonia: a prospective multicentre cohort study.}
{Pneumonia (Nathan)}
{CAPNETZ, 4,446 CAP patients with sputum; pathogen found in 1,581, cell counts in 4,014. Mixed-effects logistic models: S. pneumoniae with SO$_2$ (OR 1.75, p=0.008) and CO (1.57); E. coli with NO$_2$ (2.75); higher PM$_{2.5}$ with lower detection of the pathogen examined (OR 0.37, p=0.001). No CIs or increments in the abstract, so magnitudes are ungradable. City-level exposure over a month is coarse, many pollutant-pathogen pairs were tested, and an inverse PM$_{2.5}$ association invites a sampling or antibiotic-pretreatment explanation the abstract does not exclude. Relevance: indirect - outcome definition, not exposure metrology, limits this.}
{10.1186/s41479-026-00212-3}

\par\vskip 0pt plus 10\baselineskip\penalty-200\vskip 0pt plus -10\baselineskip
\band{Mechanistic toxicology\hfill 1 record}{Violet}

\paperentry{Sun Waite review molecular mechanisms respirable silica coal dust asbestos lung disease 2026}
{Molecular mechanisms of dust-induced lung disease: exploring the role of physicochemical composition and cellular pathways.}
{Part Fibre Toxicol}
{Review of respirable occupational dust toxicity: oxidative stress, mitochondrial and lysosomal damage, inflammasome activation, ferroptosis, necrosis and pyroptosis, macrophage/neutrophil activation and NK dysfunction driving fibrosis and tumorigenesis, with emphasis on silica and on coal dust bearing reduced iron. Covers antifibrotic and immunotherapy approaches. Narrative, no stated search or appraisal. The physicochemical angle (surface iron, crystallinity) is the useful part: it is why mass concentration alone mis-ranks hazard across dust types. Relevance: indirect - optical low-cost sensors are blind to mineralogy, so occupational dust monitoring needs composition or a source-specific mass-to-hazard factor.}
{10.1186/s12989-026-00696-3}

\par\vskip 0pt plus 10\baselineskip\penalty-200\vskip 0pt plus -10\baselineskip
\band{Exposure assessment \& modelling\hfill 9 records}{Clay}

\paperentry{Tan Shenzhen PMF-ML firework burning PM$_{2.5}$ source apportionment 2026}
{Apportioning Firework Burning and Routine Source Contributions to PM(2.5) in a Chinese Megacity Using Two PMF-Machine Learning Frameworks.}
{Environ Res}
{Three Spring Festival seasons of hourly speciation in Shenzhen; PMF resolved 11 routine sources plus firework burning. Fireworks contributed 4.7\% of PM$_{2.5}$ over the full period, 33\% during display episodes and over 50\% at peak hours. Two ML surrogates (direct and indirect) reproduced PMF contributions and flagged fireworks as dominant during episodes. Agreement with PMF is circular validation - the ML learns PMF output, so it inherits PMF's factor ambiguity and no independent tracer test is reported. Relevance: direct for networks - episodic sources at >50\% of hourly mass are exactly where a fixed humidity/composition calibration on optical sensors breaks.}
{10.1016/j.envres.2026.126121}

\paperentry{Sathish Kanpur Mumbai nine-stage size-resolved aerosol chemistry ICRP deposited dose 2026}
{Size-resolved measurements matter for chemistry and health: paired inland-coastal size-resolved aerosol chemistry and the inadequacy of bulk PM(2.5) dose estimates.}
{Environ Res}
{Paired Kanpur (inland IGP) and Mumbai (coastal) impactor sampling over three seasons. Fine mode (PM2.1) is up to 62\% of mass in Kanpur, not Mumbai; Kanpur fine aerosol is more acidic (pH 3.5-4.5 vs 4.5-5.0). ICRP 66 dosimetry: fine mode delivers 53-75\% of alveolar dose from 40-62\% of mass, and Kanpur alveolar doses are 1.3-2 times Mumbai's. The authors flag that SOC-water associations lose significance at one value per sample day, and dose depends on assumed deposition fractions. Relevance: direct - a calibration fitted at a coastal site will misweight the inland fine mode that carries the dose.}
{10.1016/j.envres.2026.126094}

\paperentry{Ayittey Accra FRM gravimetric PM$_{2.5}$ three sites e-waste 2021-2022 2026}
{Persistent urban PM$_{2.5}$ pollution in Accra, Ghana: spatiotemporal patterns, meteorological and anthropogenic drivers, and associated health risks}
{Environ Monit Assess}
{Two years (2021-2022) of FRM gravimetric PM$_{2.5}$ at Adabraka (next to the Agbogbloshie e-waste burning area), Dansoman and Legon. All sites ran about 3-4 times the WHO 24-h guideline; Adabraka was highest at 60.76 (2021) and 60.54 $\mu$g\,m$^{-3}$ (2022), with negligible interannual change but strong spatial contrast. PCA ties variability positively to temperature and negatively to RH. Hazard quotients exceeded 1 for adults and children. Sampling was intermittent (\textasciitilde{}300 samples over 3 sites and 2 years), and HQ applied to PM mass is a screening device, not an epidemiological risk. Relevance: direct - one of very few reference-grade West African series, usable to anchor low-cost and satellite calibration. \textsc{backfill}, entry 3 March 2026.}
{10.1007/s10661-026-15117-3}

\paperentry{Jaelani Surabaya Benowo landfill 15-min PM$_{2.5}$ 2020-2024 2026}
{The spatial impact assessment of landfill sites on air quality in Surabaya: a diurnal and seasonal perspective}
{Environ Sci: Atmos}
{Only a one-sentence abstract is available. Five years of 15-minute monitoring at the Benowo landfill show PM$_{2.5}$ rising from 37.17 to 50.32 $\mu$g\,m$^{-3}$ between 2020 and 2024, 7-10 times the WHO annual guideline, with exceedance peaks at 07:00 and overnight. The authors attribute the timing to boundary-layer mixing rather than landfill operations. Instrument class, completeness and any upwind comparison are not stated, so landfill attribution cannot be judged. Relevance: indirect - the diurnal peak set by mixing depth is a reminder that a single fixed node cannot separate source from meteorology without a paired background site.}
{10.1039/d6ea00048g}

\paperentry{Ku North China Plain new particle formation GEOS-Chem-APM spring campaign 2026}
{Investigations of New Particle Formation Mechanisms in North China Plain during the Spring Campaign Based on Observations and Simulations}
{Environ Sci Technol}
{One-month 2016 spring campaign in an industrial North China Plain city: 10 NPF days, each after nighttime northwesterly ventilation under dry daytime conditions. GEOS-Chem-APM reproduced 6 of 10 events; NH$_3$ was in excess; tripling H2SO4 raised nucleation but did not reproduce particle evolution, and simulated growth rates were 51\% and 75\% too low on well- and poorly-simulated days. Ten events and one season limit generality, and the campaign is a decade old. No mass or health metric. Relevance: indirect - NPF bursts raise ultrafine number with little PM$_{2.5}$ mass, the regime mass-calibrated optical nodes cannot see.}
{10.1021/acs.est.6c07540}

\paperentry{Schwarz Helmholtz letter ultrafine particle health studies forest for the trees 2026}
{Letter to the Editor: Health effect studies of ultrafine air pollution-Are we missing the forest for the trees?}
{Environ Res}
{Metadata only; a letter with no abstract, so no argument is asserted beyond the title. From the Helmholtz Munich epidemiology group (Schwarz, Pickford, Cyrys, Breitner-Busch), which runs long-term UFP number and size measurement in Augsburg. The title signals a critique of how UFP health studies are framed - plausibly metric choice (number vs surface vs size bins) or the pursuit of UFP-specific effects independent of correlated traffic pollutants. To be read when the full text is accessible. Relevance: possible - any UFP exposure standard decides what number-concentration sensors in a network should report.}
{10.1016/j.envres.2026.126091}

\paperentry{Hussain Bergin Pakistan smog episodes data-driven source attribution 2026}
{Attribution of Air Pollution Sources During Smog Episodes Across Major Urban Centers of Pakistan: A Data-Driven Approach}
{Atmos Environ}
{\textsc{metadata only.} Metadata only; no abstract in Crossref at issue time, so no finding is asserted. The title promises data-driven source attribution for smog episodes across several Pakistani cities; the author list includes Mike Bergin, whose group has built low-cost and satellite PM$_{2.5}$ work in South Asia, and Muhammad Fahim Khokhar. Carried because Punjab smog is among the most extreme PM regimes, and attribution there decides between crop burning, brick kilns and traffic. To be assessed when the abstract deposits: data sources (sensor, reference, satellite) and method. Relevance: possible, unassessed.}
{10.1016/j.atmosenv.2026.122473}

\paperentry{Shankar meteorological drivers urban air quality weather-pollution interactions 2026}
{Understanding the Meteorological Drivers of Urban Air Quality: Weather-Air Pollution Interactions}
{Atmos Environ}
{\textsc{metadata only.} Metadata only; no abstract in Crossref at issue time, so no finding is asserted. The title points to an analysis, or possibly a review, of how weather variables drive urban air quality. Carried provisionally: meteorological normalisation is a prerequisite for interpreting network trends and for separating calibration drift from real change. Geography, pollutants and design are unknown, and it may be gas-dominated, in which case it will be noted as such in the weekly. Relevance: possible, unassessed.}
{10.1016/j.atmosenv.2026.122468}

\paperentry{Yao Harbin early-spring heavy pollution episodes natural anthropogenic drivers 2026}
{Natural and anthropogenic drivers of heavy pollution episodes in early spring in Harbin, Northeast China}
{Atmos Pollut Res}
{\textsc{metadata only.} Metadata only; no abstract in Crossref at issue time, so no finding is asserted. The title sets natural against anthropogenic drivers of early-spring heavy pollution in Harbin, a period where heating-season combustion, agricultural residue burning and stagnant boundary layers overlap. Carried because cold, humid episodes are the regime where optical low-cost sensors overestimate most and where humidity correction matters. To be assessed when the abstract deposits: whether it separates open burning from heating and uses PM composition. Relevance: possible, unassessed.}
{10.1016/j.apr.2026.103263}

\par\vskip 0pt plus 10\baselineskip\penalty-200\vskip 0pt plus -10\baselineskip
\band{Sensing, forecasting \& instrumentation\hfill 4 records}{Amber}

\paperentry{Agrawal Bihar low-cost sensor network ML exposure inequity short-term mortality 2026}
{Unveiling hidden air pollution exposure and impacts with low-cost sensor network-based frameworks.}
{Nat Commun}
{Dense low-cost PM$_{2.5}$ network across Bihar, fused with ML surfaces, NFHS socioeconomic data and India-specific relative risks. Northern districts carried 1.2 times the mean exposure of southern ones; district mortality per 100,000 varied up to 1.7-fold and daily attributable mortality 3.2-fold. State-level exposure inequity was limited, intra-regional inequity substantial. The abstract gives no sensor count, calibration R2/RMSE, or uncertainty on the attributable deaths, and the mortality is borrowed risk functions applied to sensor-derived surfaces, not an observed association; network calibration error propagates straight into the burden. Relevance: direct - the clearest flagship case for low-cost networks as a burden-estimation input outside cities.}
{10.1038/s41467-026-77148-1}

\paperentry{Chen TRACER Houston micropulse lidar aerosol vertical profile boundary layer parameterisation 2026}
{Aerosol vertical distributions shaped by boundary layer dynamics in a coastal urban environment: insights from the TRACER campaign}
{Atmos Chem Phys}
{Six months of TRACER (Houston, 2022) micropulse lidar retrievals anchored to ground aerosol data. An S-shaped fitting function yields an aerosol-layer depth and a stratification-strength parameter. Concentrations drop sharply at strong capping inversions, layer depth tracks boundary-layer growth and collapse, and sea-breeze passage swaps continental for maritime aerosol. Relationships with BLH and the potential-temperature gradient are statistically significant but weak (no r given). No health endpoint. Relevance: direct for satellite-surface fusion - the AOD-to-surface PM$_{2.5}$ ratio depends on exactly this profile shape, and the parameterisation offers a physical predictor in place of BLH alone.}
{10.5194/acp-26-15543-2026}

\paperentry{Wang PM$_{2.5}$ PM$_{10}$ forecasting dual-data MWS-optimised ensemble learning 2026}
{PM$_{2.5}$ and PM$_{10}$ Concentration Forecasting Using a Dual-Data Strategy and MWS-Optimized Ensemble Learning}
{Atmos Environ}
{\textsc{metadata only.} Metadata only; no abstract in Crossref at issue time, so no finding is asserted. The title describes PM$_{2.5}$ and PM$_{10}$ forecasting from a dual-data input strategy with an ensemble learner tuned by an MWS optimiser. Carried because short-horizon forecasting is the step most networks attempt after calibration, and ensemble skill claims usually hinge on whether validation is time-blocked. To be assessed when the Elsevier abstract deposits: horizon, held-out period, skill against persistence, and whether sensor or reference data feed it. Relevance: possible, unassessed.}
{10.1016/j.atmosenv.2026.122424}

\paperentry{Rogers VUV-AMS siloxane radical mass calibration source aerosol mass spectrometry 2026}
{Siloxane-Radical Species: Chemical Characteristics and Use for Mass Calibration in Aerosol Mass Spectrometry.}
{J Phys Chem A}
{Synchrotron VUV photoionisation aerosol mass spectrometry: heating conductive silicone tubing (150-230 C) inline with a soot source deposits siloxane oligomers on particles, producing three substituted-siloxane radical series and 21 recurring peaks between 73 and 799 Da. Calibration was verified on five PAHs (178-300 Da), 2,5-dimethylfuran and ethylene-pyrolysis hydrocarbons (165-374 Da). The radical-formation mechanism on the CuO target is hypothesised, not shown, and the method is specific to one research-grade platform. Relevance: indirect - siloxane from silicone sample lines is also a known artefact; the paper doubles as a warning for any sensor or sampler plumbed with heated silicone tubing.}
{10.1021/acs.jpca.6c04791}

\par\vskip 0pt plus 10\baselineskip\penalty-200\vskip 0pt plus -10\baselineskip
\band{Occupational \& indoor\hfill 3 records}{Amber}

\paperentry{Kumar firefighter occupational exposure neurotoxicity narrative review 2026}
{Occupational exposure and neurotoxic health effects in firefighters: toxicological mechanisms and health risks.}
{Toxicol Mech Methods}
{PRISMA-informed narrative review of literature to June 2026, without meta-analysis or risk-of-bias assessment. Proposed pathways: CO and HCN hypoxia with ATP depletion acutely; fine PM and PAHs reaching the brain, NLRP3 glial activation and ferroptosis chronically; heat, fatigue and sleep loss as amplifiers. The authors state the firefighter literature is mostly cross-sectional and the mechanisms come from non-firefighter models, so causality is uncertain. Relevance: indirect - fireground PM exposure is essentially unmeasured at the individual level; wearable PM sensing on turnout gear is the obvious gap.}
{10.1080/15376516.2026.2748390}

\paperentry{Subramanian chamber size-resolved emission rates PM oxidative potential six indoor sources 2026}
{Size- and Time-Resolved Emission Rates of Particulate Matter Mass, Chemical Constituents, and Oxidative Potential from Indoor Sources}
{Environ Sci Technol}
{Chamber tests of two incenses, candle, toaster, humidifier and cigarette. Emission was mostly unimodal and concentrated below 0.25 um, where mass-normalised DTT OP was highest; composition and OPm varied with size but little with time. 37 +/- 3\% of emitted PM$_{10}$ mass and OP is deposition-capable, 45 +/- 2\% of it alveolar. For equal PM$_{10}$ mass and OP, these indoor sources give an estimated 76 +/- 10\% higher alveolar dose than ambient PM. Six sources in one chamber; the dose comparison is modelled, not measured. Relevance: direct - optical sensors undercount below \textasciitilde{}0.3 um, so indoor nodes systematically miss the fraction carrying most OP.}
{10.1021/acs.est.6c08882}

\paperentry{Gupta Maharashtra opencast coal mine PM$_{2.5}$ composition oxidative potential 2026}
{Comprehensive assessment of PM$_{2.5}$ exposures around a regional mining site in India: integrating chemical composition, oxidative potential and health risk}
{Environ Geochem Health}
{Six sites around an active opencast coal mine in eastern Maharashtra: two roadside (\textasciitilde{}400 $\mu$g\,m$^{-3}$), residential at 2 and 5 km (122 +/- 43 and 116 +/- 36 $\mu$g\,m$^{-3}$), and background. Carbonaceous species and metals were \textasciitilde{}40-50\% of mass; Cr, Ni, Fe, Mg and Zn ran 5-10-fold higher near the mine. OPvAA and OPvDTT peaked roadside and correlated only weakly with metals and organics (r=0.3-0.5); regression explained 25-40\% of OP. EC cancer risk \textasciitilde{}1e-3 rests on a risk-assessment slope, not outcomes. Season and sample counts are not in the abstract. Relevance: direct - at this composition, optical-sensor mass calibrations from urban sites will not transfer. \textsc{backfill}, entry 25 June 2026.}
{10.1007/s10653-026-03301-w}

\par\vskip 0pt plus 10\baselineskip\penalty-200\vskip 0pt plus -10\baselineskip
\band{Burden, policy \& mitigation\hfill 2 records}{Coral}

\paperentry{Laban scoping review household backup generators emissions health Africa 2026}
{Health effects associated with emissions from household generators used for backup power in Africa: a scoping review.}
{BMC Public Health}
{PubMed and Google Scholar scoping review; nine studies met inclusion, and only \textasciitilde{}2\% of screened studies addressed health outcomes directly. Generator emissions carried PM$_{2.5}$, PM$_{10}$, NOx and SO$_2$ at high levels; reported links to respiratory symptoms, CO poisoning, and possibly cancer and hypertension rest on self-report, modelling or indirect evidence. No study established a causal link. Two databases, English only. Relevance: direct - generator plumes are intermittent, near-field and household-specific, a source that only distributed indoor/outdoor low-cost sensing can capture at the needed resolution in Nigeria, Ghana or South Africa.}
{10.1186/s12889-026-28946-1}

\paperentry{US coal-fired power plant proximity Area Deprivation Index cross-sectional 2026}
{Distribution of Coal Power, Pollution, and Poverty.}
{JAMA Netw Open}
{2023 Neighborhood Atlas and census block groups: 2,783,004 people (0.83\% of the US) live within 3 miles of a coal plant, disproportionately in the most deprived ADI decile and more often American Indian/Alaska Native. The most advantaged decile had OR 0.04 (0.00-0.81) of living within 3 miles versus the most deprived - a CI spanning near-zero that signals sparse cells. Proximity stands in for exposure; no PM, metal or stack-emission data, and plant size and wind direction are ignored. Relevance: indirect - a siting map for fenceline nodes, not an exposure estimate.}
{10.1001/jamanetworkopen.2026.42335}

\par\vskip 0pt plus 10\baselineskip\penalty-200\vskip 0pt plus -10\baselineskip
\band{Other clinical endpoints\hfill 1 record}{Slate}

\paperentry{Yang US claims case-control PM$_{2.5}$ sinonasal malignant neoplasm 2026}
{Fine Particulate Matter Exposure and Sinonasal Malignant Neoplasm.}
{JAMA Otolaryngol Head Neck Surg}
{National commercial-claims case-control, 2007-2020: 3,600 SNM cases, 17,695 controls (\textasciitilde{}1:5). Per SD PM$_{2.5}$ (about 2.5 $\mu$g\,m$^{-3}$), adjusted OR 1.25 (1.05-1.49) for the 6-month mean, 1.31 (1.08-1.59) for 1 year and 1.36 (1.11-1.66) for 5 years, rising with averaging window. Exposure is assigned at metropolitan-area level, so a 1-km product is collapsed to an ecological contrast of \textasciitilde{}0.2-0.3 $\mu$g\,m$^{-3}$ between groups; smoking and deprivation are regional, not individual, and occupational wood/leather dust - the dominant SNM cause - is unmeasured. Cases are 10 years older than controls. Relevance: indirect - a rare endpoint where residence-level exposure from dense sensing would remove the MSA misclassification.}
{10.1001/jamaoto.2026.3453}

\clearpage
\section*{\sffamily\bfseries\color{Deep}Corpus register}
\vspace{-1.5mm}{\color{Amber}\rule{\linewidth}{1.1pt}}\vspace{2.5mm}

\input{table_rows}

\vspace{2mm}

\begin{keybox}
{\sffamily\bfseries\footnotesize\color{Deep}Provenance and method}\\[1.2mm]
{\fontsize{7.2}{8.8}\selectfont
Entry window \textbf{5 November 2026}, single day (Thursday). Harvester legs: \textbf{PubMed} health \textbf{12}, sensing \textbf{2}; \textbf{Europe PMC} \textbf{0}; \textbf{Crossref by ISSN} \textbf{59}; \textbf{arXiv} no in-window entry; \textbf{OpenAlex} HTTP 429 on the harvester leg. 72 unique after cross-source dedupe, 0 already carried, \textbf{72 fresh}; PubMed connector (EDAT 5 Nov, three axes) 28 PMIDs, 22 bucketed to 5 Nov, 13 appended --- \textbf{85 screened}; the 6 bucketed to 4 Nov were 3 already carried and 3 off topic. Abstracts: 25 of 26 recovered by efetch/Crossref; the Elsevier day-0 deposits stayed empty.
\textbf{Amass} BiomedCore (ingest-date filter) 15 returned: 5 already in \texttt{seen.json} or rejected, \textbf{1 admitted as backfill} (coal-mine composition), the rest previously carried or held. \textbf{Consensus} calibration sweep: 10 hits, all matched to the corpus or the rejection log except one (Gökozan, Chemosensors), held for a DOI check. OpenAlex, alphaXiv, bioRxiv/medRxiv and Wiley Scholar Gateway connectors were not available on this session; \textbf{Scite} retraction check not run (paid-plan refusal on every prior run), so no record carries a Scite check.
\textbf{25 admitted} (22 window, 3 backfill), \textbf{63 rejected} with logged reasons. \textbf{Trial watch} 4--6 November: ClinicalTrials.gov 1 hit (NCT07914010 REACT, carried); Amass TrialCore 15, relevant ones (ACTRN12626001238369, ChiCTR2600134305) already carried --- \textbf{0 new}. \textbf{5 carried on metadata alone}. Affiliations from PubMed XML (OpenAlex unavailable).
Abstracts remain the copyright of their publishers.\par}
\end{keybox}

\end{document}
